\chapter{Testing}

% intro: Boeing 747
A longs time ago, \idx{Boeing}{Boeing} was known for making products of exceptional engineering and quality. The best example of this is likely the \idx{Boeing 747}{Boeing 747}; the original \say{Jumbo Jet}. This airplane was such a leap forward that the media hailed it as a revolutionary engineering achievement. Test pilots remarked that it -- despite its massive size -- \say{landed itself}. It contained 274km of \idx{wires}{Wire} and was made up of six million physical \idx{components}{Component}. And the truely remarkable thing was that it worked \idx{reliably}{Reliability}.

% complexity compared to software
Six million components that need to work together. That certainly is a lot of complexity. But, if we look past the fact that these components still need to work together as they \idx{age}{Aging}, these are actually quite small numbers compared to what we see in software. Not that this can ever be an apples-to-apples comparison. Unlike software, airplane parts degrade over time. On the other hand, when an airplane component is degraded or faulty, it will often be easy to detect visually. Bugs in software exists precisely because they have avoided detection.

\section{Purpose}

% why we test
Given the sheer complexity of software systems and the lack of obvious indicators of faults, the need for dealing with faults is different in nature. For starters, testing -- the act of checking whether some functionality behaves as specified -- is done entirely before the product is shipped.

% test case
Testing is performed by evaluating so-called test cases. These are concrete scenarios that are well-described in terms of the intended outcome. That scenario usually covers a small subset of functionality. While it may not be packaged as such, that functionality has the interface of a function. To activate that \say{function}:
\begin{enumerate}
  \item The evaluation happens in a context with well-defined state. This must first be setup.
  \item The function needs concrete values. The combination of this and the context is the scenario.
  \item After evaluation of the function, we observe both the return value and the the context (which may have been transformed).
\end{enumerate}

% choosing test cases
The test cases are chosen in such a way that they cover the factors that influence the function and test that the effects that they produce match our perception of what is correct behavior. A well-designed test case is one whose results provide valuable insight into the quality and behavior of the code.

% test suites
As we rarely use software in as narrow situations as can be covered in a single test case, there is a need for a collection of test cases; a test suite. A test suite should be comprehensive enough to provide the evidence needed to argue that the code is probably (relatively) correct. When we run a test suite, it will tell us which test cases pass, and which fail. By looking at how that changes after a modification of our codebase, we can reason about which bugs this update fixed and whether it introduced new ones. That, of course, is only true to the granularity and correctness of our test suite.


\section{Test Types}
\subsection{Unit Testing}
\subsection{Integration Testing}


% figure: array split
\begin{figure}[tbp]
  \input{figs/dynamic/test-integration.tex}
  \caption{Integration test example.}
  \label{fig:test:integration}
\end{figure}


\subsection{System Testing}
\subsection{Acceptance Testing}
\subsection{Overview}

\section{Terminology}

% test Case

% test Suite

\section{Impossible Completeness}

% present base case
Let's define a functionality that is simple, but nontrivial. This can be concretized as a function that takes a \typename{sbyte} as parameter, adds the value \valuename{1} to it and returns the result. We could call that function \funcname{incr}.

% how to test it
To ensure \idx{correctness}{Correctness}, we must verify that all possible inputs produce the correct output. As a \typename{byte} can take on $2^8=256$ different values, we need 256 test cases. That is a lot of work, but it can be done. The job of defining these test cases requires concentration though. A mistake in a test case can easily cause an \idx{bug}{Bug} to remain undetected.

% scalability issue
But, what if the \idx{parameter}{Parameter} was a 64 bit or -- gods forbid -- a 128 bit integer? A 128 bit integer value can take on $2^{128}=3.4 \cdot 10^{38}$ different values. This is by almost all \idx{metrics}{Metric} a very large number. How large a number? There are $3.34 \cdot 10^{38}$ \idx{electrons}{Electron particle} in a cubic kilometer of \idx{water}{Water}. So, if we are able to store every test case in a single electron -- something not allowed by our current understanding of \idx{physics}{Physics} -- then we will need approx $1km^2$ of water simply to store these test cases. If we instead have two such integers as parameters, then we would need $2^{128 \cdot 2}$ electrons to store all the test cases, and this is dangerously close to the number of electrons in the known \idx{universe}{Universe}. This \idx{approach}{Approach} clearly does not \idx{scale}{Scalability}.

\section{Boundary Cases}

% leadup
As it is intractable to test all cases, we need a different approach; one that trades completeness for test execution efficiency. In short, we need to pick a limited subset of the possible inputs, and test these. But which subset do we pick? We could randomly select $n$ inputs and test these. But how do we pick one $n$ over another? If or \funcname{incr} function works correctly for an input of 12 why wouldn't it also work for a value of 13?

\subsection{Equivalence Partitions}

% figure: array split
\begin{figure}[tbp]
  \begin{center}
  \begin{tikzpicture}[]
    \newcommand{\evenY}{0mm}
    \newcommand{\oddY}{-20mm}
    \newcommand{\centerX}{0mm}
    \newcommand{\cellWidth}{12mm}
    \newcommand{\undecidedColor}{black}
    \newcommand{\beforeColor}{teal}
    \newcommand{\afterColor}{purple}
    \newcommand{\bgcolor}{white}
    \newcommand{\calcX}[2]{\centerX-#2/2*\cellWidth+#1*\cellWidth}
    
    % id, color, x, rowstride, y
    \newcommand{\cell}[5]{
      \node[cell,draw=\bgcolor!60!#2,fill=\bgcolor!60!#2] (#1) at (\calcX{#3}{#4}, #5) {};
    }
    \newcommand{\eccell}[5]{
      \node[cell,draw=#2,fill=\bgcolor!60!#2] (#1) at (\calcX{#3}{#4}, #5) {};
    }
    
    \tikzstyle{epline} = [
      overlay,
      very thick,
    ]
    \tikzstyle{ep} = [
      anchor=north,
    ]
    \tikzstyle{cell} = [
      overlay,
      very thick,
      anchor=center,
      minimum width=\cellWidth-1mm,
      minimum height=2mm,
    ]
    
    \eccell{evenI}{\beforeColor}{0}{8}{\evenY}
    \cell{evenII}{\beforeColor}{1}{8}{\evenY}
    \cell{evenIII}{\beforeColor}{2}{8}{\evenY}
    \eccell{evenIV}{\beforeColor}{3}{8}{\evenY}
    \eccell{evenV}{\afterColor}{4}{8}{\evenY}
    \cell{evenVI}{\afterColor}{5}{8}{\evenY}
    \cell{evenVII}{\afterColor}{6}{8}{\evenY}
    \eccell{evenIIX}{\afterColor}{7}{8}{\evenY}
    
    \eccell{oddI}{\beforeColor}{0}{7}{\oddY}
    \cell{oddII}{\beforeColor}{1}{7}{\oddY}
    \eccell{oddIII}{\beforeColor}{2}{7}{\oddY}
    \eccell{oddIV}{\afterColor}{3}{7}{\oddY}
    \cell{oddV}{\afterColor}{4}{7}{\oddY}
    \cell{oddVI}{\afterColor}{5}{7}{\oddY}
    \eccell{oddVII}{\afterColor}{6}{7}{\oddY}
    
    \draw[epline] ([yshift=-2mm]evenI.south west) -- ([yshift=-2mm]evenIV.south east);
    \node[ep] () at ([xshift=1.5*\cellWidth,yshift=-3mm]evenI.south) {\small \textsl{Equivalence Partition A}};
    \draw[epline] ([yshift=-2mm]evenV.south west) -- ([yshift=-2mm]evenIIX.south east);
    \node[ep] () at ([xshift=1.5*\cellWidth,yshift=-3mm]evenV.south) {\small \textsl{Equivalence Partition B}};
    
    \draw[epline] ([yshift=-2mm]oddI.south west) -- ([yshift=-2mm]oddIII.south east);
    \node[ep] () at ([xshift=1*\cellWidth,yshift=-3mm]oddI.south) {\small \textsl{Equivalence Partition A}};
    \draw[epline] ([yshift=-2mm]oddIV.south west) -- ([yshift=-2mm]oddVII.south east);
    \node[ep] () at ([xshift=1.5*\cellWidth,yshift=-3mm]oddIV.south) {\small \textsl{Equivalence Partition B}};
  \end{tikzpicture}
\end{center}

  \caption{Equivalence partitions of \funcname{zero\_first\_half} example.}
  \label{fig:test:eqpar}
\end{figure}

\section{Unit Testing}

\section{\csharp}

\section{Unit Testing}
\subsection{\csharp}
\subsection{Elixir}

\exercises{test}{Testing}
